\documentclass[12pt]{article}
\usepackage[T1]{fontenc}
\usepackage[utf8]{inputenc}
\usepackage{lmodern}
\usepackage{textcomp}
\usepackage{enumitem}
\usepackage{geometry}
\usepackage{graphicx}
\usepackage{amsmath, amssymb}
\geometry{margin=1in}
\begin{document}
\begin{center}
Machine Learning - Graduate Level Assessment 20 \
Total Marks: 40 \
Answer all questions.
\end{center}

\section*{Multiple Choice (10 marks)}
\begin{enumerate}[label=\arabic*.]
\item Consider the Bayesian network given below. How many independent parameters are needed for this Bayesian Network H -\textgreater{} U \textless{}- P \textless{}- W?
\begin{enumerate}[label=(\alph*)]
    \item 2
    \item 4
    \item 8
    \item 16
\end{enumerate}
\item Statement 1| RELUs are not monotonic, but sigmoids are monotonic. Statement 2| Neural networks trained with gradient descent with high probability converge to the global optimum.
\begin{enumerate}[label=(\alph*)]
    \item True, True
    \item False, False
    \item True, False
    \item False, True
\end{enumerate}
\item Suppose your model is overfitting. Which of the following is NOT a valid way to try and reduce the overfitting?
\begin{enumerate}[label=(\alph*)]
    \item Increase the amount of training data.
    \item Improve the optimisation algorithm being used for error minimisation.
    \item Decrease the model complexity.
    \item Reduce the noise in the training data.
\end{enumerate}
\item Which of the following best describes the joint probability distribution P(X, Y, Z) for the given Bayes net. X \textless{}- Y -\textgreater{} Z?
\begin{enumerate}[label=(\alph*)]
    \item P(X, Y, Z) = P(Y) * P(X|Y) * P(Z|Y)
    \item P(X, Y, Z) = P(X) * P(Y|X) * P(Z|Y)
    \item P(X, Y, Z) = P(Z) * P(X|Z) * P(Y|Z)
    \item P(X, Y, Z) = P(X) * P(Y) * P(Z)
\end{enumerate}
\item Which of the following best describes what discriminative approaches try to model? (w are the parameters in the model)
\begin{enumerate}[label=(\alph*)]
    \item p(y|x, w)
    \item p(y, x)
    \item p(w|x, w)
    \item None of the above
\end{enumerate}
\end{enumerate}

\section*{True / False (10 marks)}
\textit{Answer True or False}
\begin{enumerate}[label=\arabic*.]
\setcounter{enumi}{5}
\item True or False: Overfitting refers to a model that can generalize well to new data but fails to fit the training data adequately.
\item True or False: In Naive Bayes, attributes are statistically dependent of one another given the class value.
\item True or False: The variance of the Maximum A Posteriori (MAP) estimate is typically lower than that of the Maximum Likelihood Estimate (MLE).
\item True or False: An increase in P(A) while P(A, B) decreases implies that P(B|A) must also increase to maintain consistency in probabilities.
\item True or False: The cost of one gradient descent update, given the gradient vector g, is O(D), where D is the number of features.
\end{enumerate}

\section*{Long Form Response (20 marks)}
\textit{Answer each question with a clear and complete explanation.}
\begin{enumerate}[label=\arabic*.]
\setcounter{enumi}{10}
\item Discuss how the assumption of full versus diagonal class covariance matrices in a Gaussian Bayes classifier influences the trade-off between underfitting and overfitting.
\item Discuss how to generate a $10\times 5$ Gaussian matrix with $\mu=5$ and $\sigma^2=16$, and a $10\times 10$ uniform matrix from $U[-1,1)$ using PyTorch.
\end{enumerate}

\vfill
\noindent\textit{End of Paper}
\end{document}